\SourceLeaf{423}%
durch die $3$ Punkte $(1),(2),(3)$ bezeichnet werden, die also paarweise $90^{0}$ voneinander entfernt sind und zugleich den Sinn angeben sollen, in dem die Koordinaten als zunehmend betrachtet werden. Wir stellen hier einige bekannte Sätze zusammen, die wir fortwährend verwenden werden.

1) Der Winkel zwischen zwei sich schneidenden Geraden wird durch den Bogen [des Großkreises] zwischen den Punkten gemessen, die auf der Kugeloberfläche ihre Richtungen darstellen.

2) Die Lage jeder Ebene kann auf der Kugeloberfläche durch den Großkreis dargestellt werden, in dem die Kugel von einer zur ursprünglichen Ebene parallelen Ebene durch den Kugelmittelpunkt geschnitten wird.

3) Der Winkel zwischen zwei Ebenen ist gleich dem Winkel zwischen den Großkreisen, die ihre Lagen darstellen, und wird somit auch durch den Winkel zwischen den Polen der Großkreise gemessen.

4) Sind $x,y,z$ und $x',y',z'$ die Koordinaten zweier Punkte, $r$ ihr Abstand voneinander und $L$ der Punkt der Kugeloberfläche, der die Richtung der Geraden vom ersten zum zweiten Punkt darstellt, so gilt
\[
\begin{aligned}
x'&=x+r\cos(1)L\\
y'&=y+r\cos(2)L\\
z'&=z+r\cos(3)L.
\end{aligned}
\]
\par\noindent\hspace*{1.2em}5) Daraus folgt leicht, dass stets
\[
\cos(1)L^{2}+\cos(2)L^{2}+\cos(3)L^{2}=1
\]
[und] auch, wenn $L'$ irgendein anderer Punkt der Kugeloberfläche ist,
\[
\cos(1)L.\cos(1)L'+\cos(2)L.\cos(2)L'+\cos(3)L.\cos(3)L'=\cos LL'.
\]

Wir fügen noch ein anderes Theorem hinzu, das unseres Wissens sonst nirgends vorkommt und öfter von Nutzen sein kann.

Es seien $L,L',L'',L'''$ vier Punkte der Kugeloberfläche und $A$ der Winkel, den $LL'''$ und $L'L''$ an ihrem Schnittpunkt bilden. [Dann gilt]
\[
\cos LL'.\cos L''L'''-\cos LL''.\cos L'L'''=\sin LL'''.\sin L'L''.\cos A.
\]
